% ============================================================================
\section{调用约定：寄存器传参}
\label{sec:syscall-convention}
% ============================================================================

系统调用通过寄存器传递参数，这和普通的 C cdecl 函数调用（通过栈）不同。

\begin{figure}[H]
  \centering
  \begin{tikzpicture}[
    reg/.style={draw, minimum width=1.4cm, minimum height=0.7cm,
                font=\ttfamily\small, anchor=west, fill=blue!5},
    desc/.style={font=\small, anchor=west, text width=5cm},
    arrow/.style={-{Stealth[length=3pt]}, thick, blue!60!black},
  ]

  % 寄存器列
  \node[reg] (eax) at (0, 0)    {EAX};
  \node[reg] (ebx) at (0, -1.0) {EBX};
  \node[reg] (ecx) at (0, -2.0) {ECX};
  \node[reg] (edx) at (0, -3.0) {EDX};
  \node[reg] (esi) at (0, -4.0) {ESI};
  \node[reg] (edi) at (0, -5.0) {EDI};

  % 用途描述
  \node[desc] at (2.2, 0)    {syscall number（系统调用号）};
  \node[desc] at (2.2, -1.0) {arg1（第 1 个参数）};
  \node[desc] at (2.2, -2.0) {arg2（第 2 个参数）};
  \node[desc] at (2.2, -3.0) {arg3（第 3 个参数）};
  \node[desc] at (2.2, -4.0) {arg4（第 4 个参数）};
  \node[desc] at (2.2, -5.0) {arg5（第 5 个参数）};

  % 箭头
  \foreach \r in {eax, ebx, ecx, edx, esi, edi} {
    \draw[arrow] (\r.east) -- ++(0.6, 0);
  }

  % 返回值标注
  \node[reg, fill=green!10] (ret) at (0, -6.5) {EAX};
  \node[desc] at (2.2, -6.5) {返回值（负数 = $-$errno）};
  \draw[arrow, green!50!black] (ret.east) -- ++(0.6, 0);

  % 标题
  \node[font=\small\bfseries, anchor=west] at (-0.2, 1.0) {调用时};
  \node[font=\small\bfseries, anchor=west] at (-0.2, -5.8) {返回时};

  % 分隔线
  \draw[thick, dashed, gray] (-0.5, -5.9) -- (8.0, -5.9);

\end{tikzpicture}

  \caption{i386 系统调用寄存器约定（与 Linux 一致）}
  \label{fig:syscall-regs}
\end{figure}

\begin{whybox}
为什么用寄存器而不是栈传参？

\texttt{int 0x80} 触发特权级切换时，CPU 会从 TSS 读取 SS0:ESP0
并自动切换到内核栈。此时用户栈指针（SS3:ESP3）已被保存到内核栈上，
但内核\textbf{不应该信任用户栈内容}——恶意程序可能伪造栈上的数据。

寄存器在特权级切换前后由 CPU 保持不变（除了 ESP/SS 被换掉），
是最安全、最直接的传参方式。Linux i386 从 1.0 到 2.6 一直使用这套约定。
\end{whybox}

% ----------------------------------------------------------------------------
\subsection{参数快照结构体：\texttt{syscall\_regs\_t}}
% ----------------------------------------------------------------------------

后端的汇编 stub 从寄存器中提取参数后，打包成一个结构体传给 C 的 dispatch 函数。
这样 dispatch 层和具体 handler 都不依赖特定后端的寄存器布局。

\codefile{src/include/kernel/syscall.h}
\begin{lstlisting}[style=cstyle]
typedef struct syscall_regs {
    uint32_t nr;        /* EAX: 系统调用号 */
    uint32_t arg1;      /* EBX: 参数 1 */
    uint32_t arg2;      /* ECX: 参数 2 */
    uint32_t arg3;      /* EDX: 参数 3 */
    uint32_t arg4;      /* ESI: 参数 4 */
    uint32_t arg5;      /* EDI: 参数 5 */
} syscall_regs_t;
\end{lstlisting}

\begin{thinkbox}
为什么不直接把 \texttt{regs\_t}（ISR 的寄存器快照）传给 syscall handler？

\emph{提示：\texttt{regs\_t} 包含了所有通用寄存器和段寄存器（共 17 个字段），
其中大部分对系统调用无用。更重要的是，\texttt{regs\_t} 的布局和汇编 stub 的
push 顺序绑定——不同后端（int 0x80 vs sysenter）的栈布局不同。
\texttt{syscall\_regs\_t} 是后端无关的"参数传递契约"。}
\end{thinkbox}

% ----------------------------------------------------------------------------
\subsection{系统调用号定义}
\label{sec:syscall-numbers}
% ----------------------------------------------------------------------------

每个系统调用都有一个唯一的编号。用户态把编号放到 EAX，内核查表分发。
我们的编号参考 Linux i386 但不强求完全一致——只实现需要的子集。

\codefile{src/include/kernel/syscall.h}
\begin{lstlisting}[style=cstyle]
#define SYS_EXIT    1       /* 终止当前进程 */
#define SYS_WRITE   4       /* 写数据到文件描述符 */
#define SYS_BRK     45      /* 调整进程数据段末尾 */

#define SYSCALL_MAX 256     /* syscall table 容量上限 */
\end{lstlisting}

\begin{whybox}
为什么预留 0 号不使用？

如果用户态忘记设置 EAX（或 EAX 被意外清零），\texttt{syscall\_table[0] = NULL}
会让 dispatch 返回 $-$ENOSYS，而不是误调某个 handler。
这是一种防御性编程——让常见的 bug 产生可预测的错误。
\end{whybox}
